\chapter{Краткое описание основной части работы}

Основная часть состоит из шести глав: введение, обзор литературы, методология, реализация, оценка и обсуждение, заключение.

Глава~1 формулирует проблему: снапшоты содержат полный слепок памяти PVM, включая конфиденциальные данные, и хранятся без шифрования. Описаны предпосылки появления ортогональной персистентности, архитектура PhantomOS~\cite{Phantom_docs} и Genode~\cite{Genode_Foundations}, поставлены цели и задачи.

Глава~2 представляет обзор литературы по четырём направлениям. Рассматривается концепция ортогональной персистентности~\cite{PLA,Revisited,KOS,Aurora} и её реализация в PhantomOS. Описана история порта PhantomOS на Genode~\cite{Antonov_Anton} и механизм Snapper~2.0~\cite{FOSDEM_Snapper,FOSDEM_Snapper_paper}. Приводится классификация методов шифрования данных в состоянии покоя~\cite{NIST_SP800111,CryptoFS} и обосновывается выбор блочного шифрования с использованием библиотеки Tresor.

Глава~3 описывает методологию. На основе модели сценариев использования (граф состояний от включения до выключения) выбраны безопасные сценарии с аутентификацией через USB-накопитель с ключом. Сформулированы четыре требования: обязательная авторизация, хранение ключей на USB, шифрование снапшотов, гарантия записи только зашифрованных данных. Определён план реализации на базе существующего компонента \texttt{file\_vault}.

Глава~4 описывает реализацию компонента \texttt{se\_vault}. Графический интерфейс удалён, аутентификация перенесена на USB, файловая система сменена с ext2 на ext4~(\texttt{lwext4}), добавлена поддержка fsck, блочная сессия передаётся в Tresor напрямую, все блоки унифицированы до~4~КБ. В \texttt{lwext4} исправлена функция \texttt{ftruncate} и добавлена передача \texttt{sync} на нижние уровни VFS. В компонентах isomem, snapper и vfs реализована корректная цепочка завершения. Тестирование в QEMU подтвердило корректность шифрования и восстановления из снапшота. Нагрузочные тесты выявили снижение производительности в~$\approx$25~раз (первый снапшот: $\approx$2~с без шифрования и $\approx$51~с с шифрованием) из-за однопоточной реализации Tresor; инициализация ускорилась с~$\approx$20~мин до~$\approx$1~мин.

Глава~5 содержит оценку результатов. Проведён сравнительный анализ \texttt{file\_vault} и \texttt{se\_vault}; рассмотрены ограничения, связанные с непроизводственным статусом Tresor и тестированием в виртуальной среде; обсуждаются пути решения проблемы производительности.
